\documentclass[11pt,a4paper]{article}
\usepackage[T1]{fontenc}
\usepackage[utf8]{inputenc}
\usepackage{lmodern,amsmath,amssymb}
\usepackage[margin=25mm]{geometry}
\usepackage[hidelinks]{hyperref}
\hypersetup{pdftitle={The gap is bounded by the variance of the averaged spatial Polyakov loop, with the small-v},pdfauthor={navstokgap research working notes}}
\newcommand{\tightlist}{\setlength{\itemsep}{0pt}\setlength{\parskip}{0pt}}
\setlength{\emergencystretch}{3em}
\setcounter{secnumdepth}{0}
\begin{document}
\hypertarget{the-gap-is-bounded-by-the-variance-of-the-averaged-spatial-polyakov-loop-with-the-small-volume-scaling-built-in}{%
\section{The gap is bounded by the variance of the averaged spatial
Polyakov loop, with the small-volume scaling built
in}\label{the-gap-is-bounded-by-the-variance-of-the-averaged-spatial-polyakov-loop-with-the-small-volume-scaling-built-in}}

For the Kogut--Susskind Hamiltonian on the periodic lattice of side
\(L=N_sa\), the Feynman--Bijl inequality with the transverse average of
the spatial Polyakov loop in direction \(i\),
\[\bar P_i=\frac1{N_s^2}\sum_{\text{lines }\parallel i}\operatorname{Re}\operatorname{tr}\big(U_{\ell_1}\cdots U_{\ell_{N_s}}\big),\]
gives the exact bound
\[\Delta^{\rm phys}_{a,L}\ \le\ \frac{\hbar cg^2}{2a}\,\frac{\langle\sum_\ell|\nabla_\ell\bar P_i|^2\rangle_0}{\operatorname{Var}_0\bar P_i}
\ \le\ \frac{\hbar c\,g^2N}{4L\,\operatorname{Var}_0\bar P_i},\] for
every coupling and lattice size. Three consequences. In a small box at
weak coupling, where the ground state is the constant-mode state of C133
at leading order, numerator and denominator are both computable there
and the bound is \((\hbar c/L)\,g^{2/3}\) times a pure number, the
scaling of the exact zero-mode gap, so the upper side of the
small-volume corner S of \href{mass-gap-obligations-lattice.md}{the
obligations map} follows from an exact inequality plus the zero-mode
expectation values. In the infinite-volume limit, a mass gap \(m\)
forces the averaged spatial Polyakov loop to self-average at least like
\(1/L\): \(\operatorname{Var}_0\bar P_i\le\hbar cg^2N/(4mc^2L)\),
equivalently a bound on the transverse Polyakov-loop susceptibility,
\(\sum_{y}a^2\langle\operatorname{Re}P(x)\operatorname{Re}P(x+y)\rangle_c\le\hbar cg^2NL/(4mc^2)\),
which is a necessary condition for T3 expressed through one
gauge-invariant correlation. At strong coupling the averaged loop has
variance \(1/(2N_s^2)\) and the bound is empty, so the plaquette bound
of \href{lattice-gap-upper-bounds.md}{the upper-bound note} and the
present one cover complementary ends of the coupling axis. Nothing here
is promoted.

\hypertarget{the-operator-and-the-exact-bound}{%
\subsection{1. The operator and the exact
bound}\label{the-operator-and-the-exact-bound}}

Fix a direction \(i\). The lattice decomposes into \(N_s^2\) lines
parallel to \(i\), each of \(N_s\) links closing around the torus. For a
line \(\lambda\) let
\(P_\lambda=\operatorname{Re}\operatorname{tr}\prod_{\ell\in\lambda}U_\ell\),
the ordered product along the line; it is gauge invariant (a closed
loop), and \(\bar P_i=N_s^{-2}\sum_\lambda P_\lambda\). Let
\(O=\bar P_i-\langle\bar P_i\rangle_0\).

\textbf{Theorem 1.} For every \(g>0\), \(N_s\ge2\) and compact \(G\)
with the normalization of \href{mass-gap-obligations-lattice.md}{the
obligations map},
\[\Delta^{\rm phys}_{a,L}\ \le\ \frac{\hbar cg^2}{2a}\,
\frac{\big\langle\sum_\ell|\nabla_\ell\bar P_i|^2\big\rangle_0}{\operatorname{Var}_0\bar P_i}
\ \le\ \frac{\hbar c\,g^2N}{4L\,\operatorname{Var}_0\bar P_i}.\]

\emph{Proof.} \(O\) is a real gauge-invariant function of the links with
zero mean, so Proposition 1 of \href{lattice-gap-upper-bounds.md}{the
upper-bound note} applies and gives the first inequality. Each link in
direction \(i\) belongs to exactly one line and links in other
directions do not enter \(\bar P_i\), so
\(\nabla_\ell\bar P_i=N_s^{-2}\nabla_\ell P_{\lambda(\ell)}\). For a
link of \(\lambda\),
\(X_a\operatorname{tr}\prod U=\operatorname{tr}(T^aM)\) up to a phase
with \(M\) a cyclic rearrangement of the product, and the completeness
relation gives
\(\sum_a|\operatorname{tr}(T^aM)|^2=\tfrac12(N-|\operatorname{tr}M|^2/N)\le N/2\);
the real part has smaller gradient. Hence
\(\sum_\ell|\nabla_\ell\bar P_i|^2\le N_s^{-4}\cdot N_s^2\cdot N_s\cdot N/2=N/(2N_s)\),
and \(\frac{\hbar cg^2}{2a}\cdot\frac{N}{2N_s}=\frac{\hbar cg^2N}{4L}\).
\(\square\)

\hypertarget{weak-coupling-in-a-small-box}{%
\subsection{2. Weak coupling in a small
box}\label{weak-coupling-in-a-small-box}}

At weak coupling the links are near the identity and the gauge-invariant
content of a Polyakov line is the holonomy of the constant mode:
\(\prod_{\ell\in\lambda}U_\ell\approx\exp(iLa_i)\) with \(a_i=a_i^bT^b\)
the constant mode of \href{low-dimensional-mass-gap.md}{G07} Proposition
7, whose ground-state width is \(|a_i|\sim g^{2/3}/L\) in natural units.
Since \(\operatorname{tr}T^b=0\),
\[P_\lambda\approx\operatorname{Re}\operatorname{tr}e^{iLa_i}=N-\tfrac14L^2|\vec a_i|^2+O(L^4|a|^4),\qquad
\nabla_\ell P_\lambda\approx\operatorname{Im}\operatorname{tr}(T^aLa_i)=\tfrac12La_i^a ,\]
so the variance of \(\bar P_i\) is second order and the gradient first
order in the constant mode:
\[\operatorname{Var}_0\bar P_i\approx\tfrac1{16}L^4\operatorname{Var}(|\vec a_i|^2),\qquad
\Big\langle\sum_\ell|\nabla_\ell\bar P_i|^2\Big\rangle_0\approx N_s^{-4}\cdot N_s^3\cdot\tfrac14L^2\langle|\vec a_i|^2\rangle
=\frac{L^2\langle|\vec a_i|^2\rangle}{4N_s}.\] In the dimensionless
variables of C133, \(La_i=g^{2/3}\xi_i\) with \(\xi_i\) distributed by
the unit ground state \(\psi_0^{(3)}\) of
\(\tfrac12(-\Delta_{\mathbb R^9}+\sum_{i<j}|\vec\xi_i\times\vec\xi_j|^2)\):
\[\frac{\hbar cg^2}{2a}\cdot\frac{g^{4/3}\langle|\vec\xi_i|^2\rangle/(4N_s)}{g^{8/3}\operatorname{Var}(|\vec\xi_i|^2)/16}
=\frac{\hbar c}{L}\,g^{2/3}\cdot\frac{2\langle|\vec\xi_i|^2\rangle}{\operatorname{Var}(|\vec\xi_i|^2)} .\]
The bound is therefore \((\hbar c/L)g^{2/3}\) times the pure number
\(2\langle|\vec\xi_i|^2\rangle/\operatorname{Var}(|\vec\xi_i|^2)\),
finite because the unit ground state decays faster than any power along
the valleys (the confining bound of C133 gives exponential decay of
\(\psi_0\) in \(\sum_i|\vec\xi_i|\)). Consistency with the exact
zero-mode gap \(\delta_1^{(3)}g^{2/3}\hbar c/L\) requires this number to
be at least \(\delta_1^{(3)}\), which is the Feynman--Bijl inequality
applied inside the zero-mode model with the trial operator
\(|\vec\xi_i|^2\); the model version of Theorem 1 is exact and says
exactly that.

What is proved and what is not: Theorem 1 is exact. The two
approximations, replacing the ground-state expectations by their
constant-mode values, are the leading order of the small-volume
expansion (Lüscher; abstract, B78) and are not proved here. The
statement that follows from them is that the exact lattice inequality
carries the \(g^{2/3}\) scaling of the small-volume corner on its upper
side; a lower bound with that scaling remains the open weak-coupling
target of STATE.

\hypertarget{infinite-volume-a-necessary-condition-for-the-mass-gap}{%
\subsection{3. Infinite volume: a necessary condition for the mass
gap}\label{infinite-volume-a-necessary-condition-for-the-mass-gap}}

Suppose the lattice theory at fixed \(a\) and \(g\) has infinite-volume
gap \(\Delta_a=\liminf_{N_s\to\infty}\Delta_{a,L}>0\) (T2\('\)), and
write \(mc^2=\Delta_a\). Theorem 1 gives, for every \(L\),
\[\operatorname{Var}_0\bar P_i\ \le\ \frac{\hbar cg^2N}{4L\,\Delta_{a,L}} .\]
By translation invariance in the transverse plane,
\(\operatorname{Var}_0\bar P_i=N_s^{-2}\sum_{y}C_i(y)\) with
\(C_i(y)=\langle\operatorname{Re}P_\lambda(x)\operatorname{Re}P_\lambda(x+y)\rangle_c\)
the connected two-point function of the spatial Polyakov loop at
transverse separation \(y\). Hence

\textbf{Corollary 2.} If the infinite-volume gap is \(mc^2>0\), then
along any sequence of boxes on which \(\Delta_{a,L}\ge mc^2/2\),
\[\sum_{y\in\text{transverse plane}}a^2\,C_i(y)\ \le\ \frac{\hbar cg^2N}{2mc^2}\,L ,\]
and the averaged spatial Polyakov loop self-averages at least like
\(1/L\).

Reading: in a confining phase the spatial Polyakov loop of length \(L\)
has connected correlations decaying on a transverse scale \(\xi\), and
the left side is of order \(\xi^2\langle P^2\rangle_c\), so the
corollary is satisfied with room; in a phase where the loop orders or
its correlations do not decay, the left side grows like \(L^2\) and the
corollary forbids a gap. It is a necessary condition only, but it is one
written entirely in terms of a single gauge-invariant correlation and
the coupling, with no reference to the spectrum.

\hypertarget{strong-coupling-the-bound-is-empty}{%
\subsection{4. Strong coupling: the bound is
empty}\label{strong-coupling-the-bound-is-empty}}

At \(g=\infty\) the links are independent Haar variables, distinct lines
are independent, and
\(\langle|\operatorname{tr}U_1\cdots U_{N_s}|^2\rangle=1\) for a product
of independent Haar elements of \(SU(N)\) (the product is again Haar
distributed), so \(\operatorname{Var}P_\lambda=\tfrac12\) for \(N\ge3\)
and \(\operatorname{Var}\bar P_i=1/(2N_s^2)\). Theorem 1 then reads
\(\Delta\le\hbar cg^2NN_s^2/(2L)\), which grows with the lattice and
says nothing. The averaged loop is the wrong trial operator in the
disordered phase, as the plaquette is in the ordered one; the two bounds
are complementary, and neither is uniform in \(a\) along the scaling
curve.

\hypertarget{consequence-for-state}{%
\subsection{5. Consequence for STATE}\label{consequence-for-state}}

The upper side of the small-volume corner is now an exact lattice
inequality plus two zero-mode expectation values, and a mass gap has a
necessary condition in the Polyakov-loop susceptibility (Corollary 2).
The lower side, a bound \(\Delta_{a,L}\ge c\,g^{2/3}\hbar c/L\) at weak
coupling in a small box, remains the open weak-coupling target; its
natural form, suggested by the shape of Theorem 1, is an operator
inequality in which the constant-mode zero-point energy of C133 survives
the coupling to the nonzero modes.
\end{document}
